\documentclass[12pt, a4paper]{article}

% --- PREAMBLE ---
\usepackage{graphicx} % Required for inserting images
\usepackage[utf8]{inputenc}
\usepackage[T1, T2A]{fontenc} 
%\usepackage[margin=2.5cm]{geometry}
\usepackage{geometry}
\geometry{a4paper, margin=2cm}
\usepackage{csquotes}
\usepackage[english]{babel} % Russian is optional, but T2A is mandatory
\usepackage{setspace}
\setstretch{1.15} % This is roughly "1.5 line spacing"
\usepackage[skip=10pt plus 1pt, indent=0pt]{parskip}
\usepackage[T1]{fontenc} % Helps LaTeX understand where to break words
\usepackage[english]{babel} % Tells LaTeX which language's hyphenation rules to use
\usepackage{tikz}
\usepackage{amsmath, amssymb}
\usepackage{booktabs}
\usetikzlibrary{arrows.meta, positioning, shapes.geometric, shapes.symbols, calc}
\sloppy % Tells LaTeX it's okay to have slightly more space between words to avoid margin overflow

\usepackage{caption}
\usepackage{hyperref}
\usepackage{float}
\usepackage{enumitem}

\begin{document}

\input{title}

\tableofcontents

\newpage

\section{Introduction}

This chapter presents the complete mathematical and physical description of the agricultural
system studied in this thesis. The objective is to construct a high-fidelity simulation environ
ment that serves a dual purpose: first, as the virtual plant on which the Model Predictive
Controller operates, and second, as the training environment for the Reinforcement Learning
agent developed in Chapter ??. The modeling framework builds upon the agent-based model
(ABM) proposed by Lopez-Jimenez et al. [1], which couples soil water dynamics with crop
growth on a topographically heterogeneous field. However, several critical modifications and
extensions are introduced to adapt the model to the specific agronomic conditions of rice cul
tivation in Gilan Province, Iran, and to correct certain mass-balance inconsistencies identified
in the original formulation.


The chapter is organized as follows. Section 2 describes the physical field site and the digital
elevation model used to define the agent grid. Section 3 presents the complete ABM equations.
Section 4 details the mass conservation correction applied to the original water exchange for
mulation. Section 5 describes the climate data sources and the two contrasting scenarios used
throughout this thesis. Section 6 specifies the soil and crop parameters. Section 7 derives the
agronomically grounded water budget. Finally, Section 8 presents the ABM validation results,
including cross-validation against satellite soil wetness observations.

\section{Study Site Description}

The study site is a 6-hectare agricultural field located in Gilan Province, Iran (38.298$^{\circ}$N, 48.847$^{\circ}$E), in the humid subtropical coastal zone south of the Caspian Sea. Gilan is Iran's primary rice-producing province, accounting for over 238,000 hectares of paddy fields~[1]. The field exhibits a bowl-shaped topography with elevations ranging from 74~m to 181~m above sea level. A digital elevation model (DEM) was obtained from USGS Earth Explorer and processed in QGIS to extract a $10 \times 13$ grid of 130 computational agents, each representing approximately 462~m$^{2}$.

\subsection{Topological Graph Construction}

The DEM was used to construct a directed graph encoding water flow topology. For each agent~$n$, the normalized elevation was computed as $\gamma_n = (z_n - z_{\min})/(z_{\max} - z_{\min})$. Each agent sends lateral water flow to all eight-connected neighbors with strictly lower normalized elevation. The number of lower neighbors $N_r(n)$ determines the agent's drainage capacity. Agents with $N_r = 0$ are classified as \emph{sink agents}---local minima where water accumulates. On this bowl-shaped field, sink agents are concentrated near the center and represent areas at risk of waterlogging.


\section{Climate Data}

\subsection{Data Source and Period}

Daily meteorological data were obtained from the NASA POWER database for 2000--2025 at coordinates (38.30$^{\circ}$N, 48.85$^{\circ}$E), providing satellite-derived and reanalysis-based parameters at $\sim$50$\times$60~km resolution. Thirteen variables were extracted: T2M, T2M\_MAX, T2M\_MIN, PRECTOTCORR, RH2M, WS2M, PS, T2MDEW, ALLSKY\_SFC\_SW\_DWN, GWETROOT, and GWETTOP.

\subsection{Preprocessing}

The dataset was filtered to April--October for each year, yielding 5,564 daily records. No missing values were found. One rainfall anomaly was identified (April~7, 2022: 375~mm) and replaced with the calendar-day average (1.2~mm). Two reference evapotranspiration values were computed for each day (Section~\ref{sec:et0}).

\subsection{Reference Evapotranspiration}
\label{sec:et0}

Two methods were compared:

\textbf{Hargreaves equation}~[2]: $\text{ET}_0^{H} = 0.0023 \cdot (T_{\text{mean}} + 17.8) \cdot (T_{\text{max}} - T_{\text{min}})^{0.5} \cdot R_s$

\textbf{FAO Penman-Monteith}~[2]:
\begin{equation}
\text{ET}_0^{PM} = \frac{0.408 \,\Delta\,(R_n - G) + \gamma \,\frac{900}{T+273}\, u_2\, (e_s - e_a)}{\Delta + \gamma\,(1 + 0.34\, u_2)}
\end{equation}

Over the 25-year rice season, Hargreaves overestimated ET$_0$ by 34\% (6.78 vs 5.02~mm/day, $r = 0.965$). This is due to Gilan's high humidity ($\sim$64\%) and low wind ($\sim$1.40~m/s), which suppress evapotranspiration but are not captured by Hargreaves. Penman-Monteith was selected as the standard~[2].

\begin{figure}[H]
\centering
\includegraphics[width=\textwidth]{figures/fig_et0.png}
\caption{Comparison of Hargreaves and Penman-Monteith ET$_0$. Left: seasonal pattern. Right: scatter plot with regression.}
\label{fig:et0}
\end{figure}

\subsection{25-Year Climatology}

Mean temperature ranges from 12$^{\circ}$C in early April to 26$^{\circ}$C in mid-July. Precipitation shows a dry mid-season (June--August, 0.5--1.5~mm/day) flanked by wetter months. Solar radiation peaks at 26~MJ/m$^{2}$/day in June. Relative humidity remains 60--75\% throughout---a defining feature of Gilan's Caspian climate. Additional climate figures are provided in Appendix~A.

\begin{figure}[H]
\centering
\includegraphics[width=\textwidth]{figures/fig_temperature.png}
\caption{25-year daily average temperature (April--October, 2000--2025).}
\label{fig:temperature}
\end{figure}

\begin{figure}[H]
\centering
\includegraphics[width=\textwidth]{figures/fig_rainfall.png}
\caption{25-year daily average rainfall (April--October, 2000--2025).}
\label{fig:rainfall}
\end{figure}


\section{Crop Parameterization}

Two crops were selected: rice (Hashemi cultivar) and tobacco (\emph{Nicotiana tabacum}). Both are commercially significant in Gilan~[1,~3]. Their contrasting characteristics---high $K_c$/shallow roots/drought-sensitive rice versus moderate $K_c$/deep roots/drought-tolerant tobacco---enable evaluation of control algorithms under different agronomic conditions.

\subsection{Rice (Hashemi Cultivar)}

\subsubsection{Base Temperature and Thermal Time}

The base temperature $T_{\text{base}} = 10^{\circ}$C follows the standard GDD convention~[4,~5] and the baseline model~[6]. The AquaCrop convention of 8$^{\circ}$C~[20] produces proportionally higher GDD values without changing maturity dates.

The maturity threshold $\theta_{18} = 1250$~$^{\circ}$C$\cdot$day was calibrated from the 25-year dataset and validated against Sadidi Shal et al.\ (2021), who measured 1216--1352~GDD for Hashemi rice in Gilan~[7]. With a 20-day nursery contributing $\sim$210~GDD (+5$^{\circ}$C greenhouse boost), 20/26 years reach maturity by August~20.

\subsubsection{Season Window}

Transplanting: May~20 (DOY~141), when 7-day rolling $T_{\text{mean}} \geq 18^{\circ}$C and $T_{\min} \geq 12^{\circ}$C~[8,~9].

\begin{figure}[H]
\centering
\includegraphics[width=\textwidth]{figures/fig_temperature_may_years.png}
\caption{25-year average May temperature with transplanting thresholds.}
\label{fig:temp_may}
\end{figure}

Harvest: August~20 (DOY~233). WMO R20mm events ($\geq$20~mm/day)~[10] increase sharply after this date, causing lodging~[10] and harvest disruption~[11].

\begin{figure}[H]
\centering
\includegraphics[width=\textwidth]{figures/fig_rain_extremes_august.png}
\caption{Extreme rainfall events ($\geq$20~mm/day) per day of August, 2000--2025.}
\label{fig:rain_august}
\end{figure}

Field season: 93 days. Nursery provides $x_{2,\text{init}} = 210$~GDD. Other parameters: $K_c = 1.15$, $p = 0.20$~[2], $\theta_5 = 400$~mm, $\theta_{13} = 0.65$~g/MJ~[12,~13], HI = 0.42~[14]. Expected yield: 3500--4000~kg/ha (optimized) or $\sim$2000~kg/ha (traditional flooding, consistent with local farmer reports).

\subsection{Tobacco (\emph{Nicotiana tabacum})}

\subsubsection{Base Temperature and Thermal Time}

$T_{\text{base}} = 10^{\circ}$C, within the 10--13$^{\circ}$C restriction range~[15]. $\theta_{18} = 1200$~GDD (field only), calibrated to match the 104-day season. Duration consistent with FAO (90--120 days)~[16] and Britannica (100--130 days)~[17]. $x_{2,\text{init}} = 0$ (nursery conditions unknown---seedlings delivered by company).

\subsubsection{Season Window}

Transplanting: May~25 (DOY~146). Harvest: September~5 (DOY~249). 104 field days.

\subsubsection{Other Parameters}

$\theta_5 = 700$~mm (FAO: 500--1000~mm)~[16]. $K_c = 0.90$, $p = 0.50$~[16]. $\theta_9 = 30^{\circ}$C (FAO optimal ceiling)~[15,~16]. $\theta_{13} = 0.30$~g/MJ~[18], HI = 0.55. Expected yield: 2000--2500~kg/ha~[16].

\subsection{Parameter Summary}

\begin{table}[H]
\centering
\caption{Complete crop and soil parameters.}
\label{tab:params}
\small
\begin{tabular}{llccl}
\toprule
\textbf{Parameter} & \textbf{Symbol} & \textbf{Rice} & \textbf{Tobacco} & \textbf{Source} \\
\midrule
\multicolumn{5}{l}{\emph{Shared soil parameters (silty loam)}} \\
Water uptake coeff. & $\theta_1$ & 0.096 & 0.096 & [6] \\
Wilting point (vol.) & $\theta_2$ & 0.15 & 0.15 & [2] \\
SCS threshold & $\theta_3$ & 10 mm & 10 mm & [6] \\
Drainage coeff. & $\theta_4$ & 0.05 & 0.05 & [6] \\
Field capacity (vol.) & $\theta_6$ & 0.35 & 0.35 & [2] \\
\midrule
\multicolumn{5}{l}{\emph{Crop-specific parameters}} \\
Root depth & $\theta_5$ & 400 mm & 700 mm & [2], [16] \\
Base temperature & $\theta_7$ & 10$^{\circ}$C & 10$^{\circ}$C & [5], [15] \\
Heat stress onset & $\theta_9$ & 35$^{\circ}$C & 30$^{\circ}$C & [16] \\
Extreme heat & $\theta_{10}$ & 42$^{\circ}$C & 38$^{\circ}$C & -- \\
RUE & $\theta_{13}$ & 0.65 & 0.30 & [12], [18] \\
Drought sensitivity & $\theta_{14}$ & 0.80 & 0.60 & [6] \\
Maturity GDD & $\theta_{18}$ & 1250 & 1200 & [7], calib. \\
Max interception & $\theta_{19}$ & 0.95 & 0.90 & [6] \\
50\% intercept. GDD & $\theta_{20}$ & 417 & 450 & [19] \\
Crop coefficient & $K_c$ & 1.15 & 0.90 & [2], [16] \\
Depletion fraction & $p$ & 0.20 & 0.50 & [2], [16] \\
Harvest index & HI & 0.42 & 0.55 & [14], [16] \\
Nursery GDD & $x_{2,\text{init}}$ & 210 & 0 & Calib. \\
Initial biomass & $x_{4,\text{init}}$ & 60 g/m$^2$ & 8 g/m$^2$ & Est. \\
Season & -- & DOY 141--233 & DOY 146--249 & This study \\
Field days & -- & 93 & 104 & -- \\
FC & -- & 140 mm & 245 mm & $\theta_6 \times \theta_5$ \\
WP & -- & 60 mm & 105 mm & $\theta_2 \times \theta_5$ \\
TAW & -- & 80 mm & 140 mm & FC $-$ WP \\
RAW & -- & 16 mm & 70 mm & $p \times$ TAW \\
\bottomrule
\end{tabular}
\end{table}


\section{Water Budget and Scenario Selection}

\subsection{Water Budget Formulation}

Crop water demand: $\text{ET}_c = \overline{\text{ET}_0^{PM}} \cdot K_c \cdot N_{\text{days}}$. Full irrigation need: $I_{\text{full}} = \max(\text{ET}_c - \overline{P}, 0)$, using 25-year average rainfall (Method~B).

\begin{table}[H]
\centering
\caption{Water budget summary.}
\label{tab:budget}
\begin{tabular}{lcc}
\toprule
& \textbf{Rice} & \textbf{Tobacco} \\
\midrule
Season duration & 93 days & 104 days \\
Avg ET$_0$ (PM) & 5.02 mm/day & 4.95 mm/day \\
Crop demand & 537.0 mm & 463.4 mm \\
Avg season rainfall & 53.3 mm & 74.5 mm \\
Full irrigation need & 483.7 mm & 389.0 mm \\
Budget at 70\% & 338.6 mm & 272.3 mm \\
Budget at 50\% & 241.8 mm & 194.5 mm \\
\bottomrule
\end{tabular}
\end{table}

Rice demand (537~mm) is consistent with Hashemi field measurements: 534~mm (SRI) to 632~mm (traditional flooding)~[14]. Tobacco demand (463~mm) falls within the FAO range of 400--600~mm~[16]. Despite a longer season, tobacco needs less irrigation due to lower $K_c$ and higher rainfall.

\subsection{Scenario Selection}

\begin{table}[H]
\centering
\caption{Selected climate scenarios.}
\begin{tabular}{lccc}
\toprule
\textbf{Scenario} & \textbf{Year} & \textbf{Rice (mm)} & \textbf{Tobacco (mm)} \\
\midrule
Dry & 2022 & 39.7 & 54.0 \\
Moderate & 2020 & 42.1 & 73.3 \\
Wet & 2024 & 176.8 & 163.2 \\
\bottomrule
\end{tabular}
\end{table}


\section{Agent-Based Model}

The ABM follows L\'{o}pez-Jim\'{e}nez et al.~[6]. Each of $N = 130$ agents maintains: $x_1$ (soil water, mm), $x_2$ (thermal time, $^{\circ}$C$\cdot$day), $x_3$ (maturity stress), $x_4$ (biomass, g/m$^2$).

\subsection{State Equations}

\textbf{Water balance:}
\begin{equation}
x_1^n(k\!+\!1) = x_1^n(k) - \varphi_1^n - \varphi_2^n - \varphi_3^n - \varphi_{4,\text{out}}^n + \varphi_{4,\text{in}}^n + \varphi_{2,\text{in}}^n + P(k) + u^n(k)
\label{eq:water}
\end{equation}

\textbf{Thermal time:} $x_2^n(k+1) = x_2^n(k) + \max(T_{\text{mean}} - \theta_7, 0)$

\textbf{Maturity stress:} $x_3^n(k+1) = x_3^n(k) + \theta_{11}(1-h_2) + \theta_{12}(1-h_3^n)$

\textbf{Biomass:}
\begin{equation}
x_4^n(k\!+\!1) = x_4^n(k) + \theta_{13} \cdot h_3^n \cdot h_6^n \cdot h_7 \cdot g(x_2^n) \cdot R_s(k)
\label{eq:biomass}
\end{equation}

\subsection{Water Fluxes}

$\varphi_1^n = \min(\max(\theta_1(x_1^n - \theta_2\theta_5), 0), \text{ET}_0)$ (transpiration)

$\varphi_2 = (P-\theta_3)^2/(P+4\theta_3)$ if $P > \theta_3$, else 0 (SCS runoff)

$\varphi_3^n = \theta_4 \cdot \max(x_1^n + P - \varphi_2 - \theta_6\theta_5, 0)$ (deep drainage)

$\varphi_{4,\text{out}}^n = \max(x_1^n - \theta_6\theta_5, 0)$; $\varphi_{4,\text{in}}^m = \sum_{n \in \mathcal{U}(m)} \varphi_{4,\text{out}}^n / N_r(n)$ (lateral flow)

\subsection{Stress and Growth Functions}

$h_1 = \max(T_{\text{mean}} - \theta_7, 0)$; $h_2$: linear between $\theta_9$ (no stress) and $\theta_{10}$ (full stress); $h_3^n = 1 - \theta_{14} \cdot \max(1 - \varphi_1^n/\text{ET}_0, 0)$ (drought); $h_6^n = \text{clip}(1 - (x_1^n - \theta_6\theta_5)/(\theta_6\theta_5), 0, 1)$ (waterlogging); $h_7 = 1$ if $T_{\text{mean}} > \theta_7$, else 0.

\textbf{Growth function:}
\begin{equation}
g(x_2^n) = \begin{cases}
\theta_{19}/(1 + e^{-0.01(x_2^n - \theta_{20})}) & x_2^n \leq \theta_{18}/2 \\
\theta_{19}/(1 + e^{0.01(x_2^n + x_3^n - \theta_{18})}) & x_2^n > \theta_{18}/2
\end{cases}
\end{equation}


\section{Model Corrections and Extensions}

Three corrections to the baseline~[6] were implemented.

\subsection{Mass Conservation}

The original water balance includes $\varphi_{4,\text{in}}$ but omits $\varphi_{4,\text{out}}$, violating mass conservation. The corrected Eq.~\ref{eq:water} includes $-\varphi_{4,\text{out}}^n$.

\subsection{Drought Stress on Biomass}

The original biomass equation lacks $h_3$, so drought does not reduce biomass. The corrected Eq.~\ref{eq:biomass} includes $h_3^n$, reducing biomass by up to $\theta_{14} \times 100\%$ under maximum drought (80\% for rice, 60\% for tobacco). This is essential for irrigation optimization: without it, the controller has no incentive to irrigate.

\subsection{Surface Runoff Redistribution}
\label{sec:runoff}

Originally, $\varphi_2$ leaves the system. A simple redistribution was implemented: each agent's runoff flows to lower neighbors with immediate infiltration ($\varphi_{2,\text{in}}^n$). A cascading alternative (top-to-bottom SCS re-application) was also tested.

\begin{table}[H]
\centering
\caption{Runoff redistribution comparison, no irrigation, wet scenario (2024).}
\label{tab:runoff}
\begin{tabular}{lccc}
\toprule
\textbf{Mode} & \textbf{Yield (kg/ha)} & \textbf{$x_1$ sinks (mm)} & \textbf{$x_1$ hills (mm)} \\
\midrule
None (original) & 2135 & 71.2 & 71.2 \\
Simple (A) & 2273 & 74.8 & 72.0 \\
Cascade (B) & 2278 & 79.5 & 72.0 \\
\bottomrule
\end{tabular}
\end{table}

Simple redistribution captures the first-order effect (6.5\% yield improvement) with negligible difference from cascade (0.2\%). Simple was adopted.

\begin{figure}[H]
\centering
\includegraphics[width=\textwidth]{figures/fig_runoff_comparison.png}
\caption{No-irrigation runoff mode comparison, wet scenario (2024).}
\label{fig:runoff}
\end{figure}

\subsection{Biomass Calibration}

$\theta_{13}$ was calibrated so $x_4$ represents g/m$^2$ dry matter. Yield: $Y = x_4 \times \text{HI} \times 10$ (kg/ha). Rice: validated against $\sim$2000~kg/ha (farmer reports) and 3500--4000~kg/ha (optimized)~[14]. Initial conditions: rice $x_{2,\text{init}} = 210$, $x_{4,\text{init}} = 60$; tobacco $x_{2,\text{init}} = 0$, $x_{4,\text{init}} = 8$.


\section{Model Validation}

\subsection{GWETROOT Cross-Validation}

ABM soil water was compared to NASA MERRA-2 GWETROOT. Both normalized: ABM as $(x_1 - \text{WP})/(\text{FC} - \text{WP})$; GWETROOT as native 0--1 fraction. Pearson correlation assessed temporal pattern agreement.

\begin{table}[H]
\centering
\caption{Cross-validation: ABM vs NASA GWETROOT.}
\label{tab:crossval}
\begin{tabular}{lccc}
\toprule
\textbf{Scenario} & \textbf{Year} & \textbf{Pearson $r$} & \textbf{$p$-value} \\
\midrule
Dry & 2022 & 0.738 & $3.1 \times 10^{-17}$ \\
Moderate & 2020 & 0.534 & $3.4 \times 10^{-8}$ \\
Wet & 2024 & 0.337 & $9.6 \times 10^{-4}$ \\
\bottomrule
\end{tabular}
\end{table}

\begin{figure}[H]
\centering
\includegraphics[width=\textwidth]{figures/fig_gwetroot_timeseries.png}
\caption{Normalized soil moisture: ABM (blue) vs NASA GWETROOT (brown).}
\label{fig:gwetroot}
\end{figure}

The dry-year $r = 0.74$ ($p < 0.001$) confirms the ABM captures wetting/drying dynamics. Decreasing correlation with increasing rainfall is expected (reduced signal contrast in wet conditions). The absolute level difference (ABM $\sim$0.08 vs GWETROOT $\sim$0.46) reflects: (1) NASA's deeper root zone ($\sim$1000~mm vs 400~mm); (2) GWETROOT observes irrigated fields while the ABM simulates rain-fed conditions; (3) different soil parameterizations.


\section*{References}

\begin{enumerate}[label={[\arabic*]}, leftmargin=*]
\item SurfIran, ``Gilan's rice fields,'' 2025. Available: \url{https://surfiran.com/mag/rice-transplantation-in-gilan-province/}
\item R.~G.~Allen et al., ``Crop evapotranspiration,'' FAO Irrig. Drain. Paper 56, 1998.
\item ``Tobacco,'' \emph{Encyclopaedia Iranica}. Available: \url{https://www.iranicaonline.org/articles/tobacco/}
\item G.~S.~McMaster and W.~W.~Wilhelm, ``Growing degree-days,'' \emph{Agric.\ For.\ Meteorol.}, vol.~87, pp.~291--300, 1997.
\item S.~Laohasiriwong et al., ``High temperature alters phenology in rice,'' \emph{Plants}, vol.~12, no.~3, p.~666, 2023.
\item A.~L\'{o}pez-Jim\'{e}nez et al., ``Agent-based MPC of soil-crop irrigation,'' \emph{Control Eng.\ Pract.}, vol.~150, 2024.
\item S.~M.~T.~Sadidi Shal et al., ``GDD of Hashemi rice in Guilan,'' \emph{J.\ Climate Res.}, vol.~1400(45), pp.~143--152, 2021.
\item P.~Paredes et al., ``Base and upper temperature thresholds for GDD,'' \emph{Agric.\ Water Manag.}, vol.~319, 2025.
\item R.~Faghani et al., ``Planting date and seedling age effects on rice yield,'' \emph{Afr.\ J.\ Biotechnol.}, 2011.
\item A.~A.~Sarma et al., ``Adaptation of wetland rice to extreme weather,'' 2018.
\item S.~Nayak et al., ``Extreme rainfall indices and rice productivity,'' 2011.
\item S.~R.~Patel et al., ``RUE in irrigated rice,'' \emph{Paddy Water Environ.}, vol.~11, pp.~477--486, 2013.
\item V.~O.~Sadras and R.~Trenchero, ``Radiation use efficiency,'' in \emph{Crop Physiology}, Elsevier, 2021.
\item N.~Jalali et al., ``Water requirement of Hashemi rice,'' \emph{Water Res.\ Agric.}, 2021.
\item Y.~Y.~Li et al., ``Cold stress effects on tobacco,'' \emph{BMC Plant Biol.}, vol.~21, p.~131, 2021.
\item FAO, ``Crop information: Tobacco.'' Available: \url{https://www.fao.org/land-water/databases-and-software/crop-information/tobacco/en/}
\item ``Tobacco,'' \emph{Encyclopaedia Britannica}, 1999.
\item R.~Ferrara, ``RUE in flue-cured tobacco,'' \emph{Field Crops Res.}, 2002.
\item Y.~Zhang et al., ``LAI simulation of tobacco,'' \emph{Chin.\ J.\ Eco-Agric.}, 2017.
\item R.~S.~Raoufi et al., ``Climate changes and rice cultivation,'' \emph{Int.\ J.\ Plant Prod.}, 2025.
\item TasteIran, ``Iran rice harvesting.'' Available: \url{https://www.tasteiran.net/stories/10/iran-rice-harvesting}
\end{enumerate}


\appendix
\section{Supplementary Climate Figures}

\begin{figure}[H]
\centering
\includegraphics[width=\textwidth]{figures/fig_radiation.png}
\caption{25-year daily average solar radiation, April--October.}
\end{figure}

\begin{figure}[H]
\centering
\includegraphics[width=\textwidth]{figures/fig_humidity.png}
\caption{25-year daily average relative humidity, April--October.}
\end{figure}

\begin{figure}[H]
\centering
\includegraphics[width=\textwidth]{figures/fig_wind.png}
\caption{25-year daily average wind speed, April--October.}
\end{figure}

\begin{figure}[H]
\centering
\includegraphics[width=\textwidth]{figures/fig_pressure.png}
\caption{25-year daily average surface pressure, April--October.}
\end{figure}

\begin{figure}[H]
\centering
\includegraphics[width=\textwidth]{figures/fig_soil_moisture.png}
\caption{25-year daily average NASA MERRA-2 soil wetness, April--October.}
\end{figure}

\begin{figure}[H]
\centering
\includegraphics[width=\textwidth]{figures/fig_temperature_may_years.png}
\caption{Daily temperature in May for 2021--2025.}
\end{figure}

\begin{figure}[H]
\centering
\includegraphics[width=\textwidth]{figures/fig_runoff_spatial.png}
\caption{Spatial distribution of $x_1$ at harvest, three runoff modes, wet scenario (2024).}
\end{figure}

\begin{figure}[H]
\centering
\includegraphics[width=\textwidth]{figures/fig_gwetroot_scatter.png}
\caption{Scatter plots: ABM vs NASA GWETROOT for three scenarios.}
\end{figure}

\end{document}
